\section{结论}
\label{sec:conclusion}

本文构建 AquaDiag 基准，将低资源水产诊断解耦为"病灶识别（确定性）"与"病种诊断（不确定性，含合规+拒答）"，并发现三个 VLM 在该任务上存在"追问不足、拒答不足、违规开药"的系统性失败。ProDiag 以"视觉显式化 + 规则表 + 红线过滤"的轻量方法，不训练即实现诊断 74.1\% + 合规 0\% + 追问 95.7\% + 拒答 100\%，证明可审计的领域规则结构是低资源专业诊断的可行路线。专家标注进一步揭示病灶识别与病种诊断的天然分离，为后续"训练式追问/拒答对齐"奠定基础。未来工作包括：扩大病种与图量、将病灶识别评测升级为专家体征描述匹配、以及用强化学习把规则行为蒸馏进可端到端训练的模型。
